% Synthetic data generation — methodology section.
% Describes experiments/synthetic_identifiable.yaml; generator in
% eval/synthetic/synthetic_dataset.py (PseudoMRIRenderer).

\subsection{Synthetic paired-contrast data}
\label{sec:synthetic-data}

Since real paired MRI has no ground-truth latents, we evaluate on a procedural
generator that renders a T1-like and a FLAIR-like $64^3$ volume of the same
subject. Each subject has a content vector $\mathbf{c}\in[-1,1]^{9}$ shared by
both views and a style vector $\mathbf{s}^{(v)}\in\mathbb{R}^{3}$ drawn
independently for each view.

\paragraph{Content.}
The nine factors are brain size, ventricle size, lesion position (3D),
cortical thickness, temporal-lobe atrophy, left--right asymmetry and sulcal
depth. They are sampled from a nonlinear structural causal model over a random
DAG, so the factors are dependent, and are then mapped monotonically to
$[-1,1]$. Each factor deforms a tissue-label map of white matter, grey matter,
ventricles and a white-matter lesion. Local factors are constructed so that
they do not change total brain volume, which keeps them distinguishable from
brain size.

\paragraph{Style.}
Each view assigns intensities to the tissue labels with a modality-specific
lookup table (e.g.\ white matter brighter than grey matter in T1, and the
reverse in FLAIR; the lesion is hypointense in T1 and hyperintense in FLAIR).
The style vector sets a global intensity gain, an offset and the noise level.
A smooth bias field and Rician noise are then applied. Since style is redrawn
for each view, it differs within every positive pair, whereas content is
identical.
